% Размер бумаги А4, шрифт 12 пунктов
\documentclass[a4paper,12pt]{report}
% Перекодировка шрифтов
\usepackage[T2A]{fontenc}
% Кодировка: utf8
\usepackage[utf8]{inputenc}
% Используем русский и английский языки с переносами
\usepackage[english,russian]{babel}
% Подключаем нужные пакеты расширений
\usepackage{amssymb,amsfonts,amsmath,mathtext,cite,enumerate,float}

% Вставка рисунков
\usepackage{graphicx}
% Путь к рисункам
\graphicspath{{../pictures/}}

% Подсчёт объектов
\usepackage{totcount}
\regtotcounter{page}
\regtotcounter{figure}
\regtotcounter{table}

% Изменение размеров шрифтов
\usepackage{scrextend}
% Размер по умолчанию — 14 пунктов, интервал полуторный (14*1.25)
\changefontsizes[17.5pt]{14pt}
% Отступ первой строки параграфа после начала раздела
\usepackage{indentfirst}
% Отступ начала параграфа
\setlength{\parindent}{1cm}

% Отступы слева для списков
% Первый уровень выровнен на отступ нового абзаца, 
% остальные -- на начало элемента более высокого уровня.
\usepackage{enumitem}
\setlist[itemize,1]{leftmargin=1.55cm, label=---}
\setlist[enumerate,1]{leftmargin=1.55cm}
\setlist[itemize,2,3]{leftmargin=*, label=---}
\setlist[enumerate,2,3]{leftmargin=*}

% Заменяем библиографию с квадратных скобок на точку
\makeatletter
\renewcommand{\@biblabel}[1]{#1.}
\makeatother

% Геометрия страницы
\usepackage{geometry}
% Поля
\geometry{left=3cm}
\geometry{right=1cm}
\geometry{top=2cm}
\geometry{bottom=2cm}

% Меняем везде перечисления на цифра.цифра
\renewcommand{\theenumi}{\arabic{enumi})}
\renewcommand{\labelenumi}{\arabic{enumi})}
\renewcommand{\theenumii}{\arabic{enumii})}
\renewcommand{\labelenumii}{\arabic{enumi}.\arabic{enumii})}
\renewcommand{\theenumiii}{\arabic{enumiii})}
\renewcommand{\labelenumiii}{\arabic{enumi}.\arabic{enumii}.\arabic{enumiii})}

% Изменяем \part и \chapter из report.cls чтобы они не выводили нумерованные заголовки
% Рецепт отсюда: http://tex.stackexchange.com/questions/70621/part-and-chapter-commands-output-both-numbered-and-named-headings/70624
\makeatletter
    \def\@part[#1]#2{%
        \ifnum \c@secnumdepth >-2\relax
        \refstepcounter{part}%
        \addcontentsline{toc}{part}{#1}%
        \else
        \addcontentsline{toc}{part}{#1}%
        \fi
        \markboth{}{}%
        {
            \centering
            \interlinepenalty \@M
            \normalfont
            \Huge \bfseries #2\par}%
        \@endpart}
    \def\@chapter[#1]#2{\ifnum \c@secnumdepth >\m@ne
        \refstepcounter{chapter}%
        \typeout{\@chapapp\space\thechapter.}%
        \addcontentsline{toc}{chapter}%
        {#1}%
        \else
        \addcontentsline{toc}{chapter}{#1}%
        \fi
        \chaptermark{#1}%
        \addtocontents{lof}{\protect\addvspace{10\p@}}%
        \addtocontents{lot}{\protect\addvspace{10\p@}}%
        \if@twocolumn
        \@topnewpage[\@makechapterhead{#2}]%
        \else
        \@makechapterhead{#2}%
        \@afterheading
        \fi}
    \def\@makechapterhead#1{%
        \vspace*{50\p@}%
        {\parindent \z@ \raggedright \normalfont
            \interlinepenalty\@M
            \Huge \bfseries #1\par\nobreak
            \vskip 40\p@
    }}
    \renewcommand{\thesection}{\arabic{section}} % Удаляем \thechapter из \thesection
\makeatother

% Нумерация подразделов
\setcounter{secnumdepth}{3}
% Включение подразделов в оглавление
\setcounter{tocdepth}{3}
% Для веб-адресов в библиографии
\usepackage{url}
% Для подчёркиваний в полях библиографических записей
\usepackage[strings]{underscore}
% Библиография в алфавитном порядке с цифровой нумерацией
\bibliographystyle{ugost2008}
% Определяем названия месяцев вручную
\def\bbljan{Январь}
\def\bblfeb{Февраль}
\def\bblmar{Март}
\def\bblapr{Апрель}
\def\bblmay{Май}
\def\bbljun{Июнь}
\def\bbljul{Июль}
\def\bblaug{Август}
\def\bblsep{Сентябрь}
\def\bbloct{Октябрь}
\def\bblnov{Ноябрь}
\def\bbldec{Декабрь}

% Выключаем переносы
\makeatletter\chardef\l@nohyphenation=255 \makeatother
\usepackage[english=nohyphenation,russian=nohyphenation]{hyphsubst}
% Не хотим строки, которые влезают на поля
\sloppy

% !TeX spellcheck = ru_RU
\begin{document}
% Титульный лист
\begin{titlepage}
    \newgeometry{top=2cm,bottom=2cm,left=2cm,right=2cm}
    \newpage
    
    \begin{center}
        Федеральное агенство по образованию Российской Федерации \\
        Московский Государственный Технический Университет \\*
        имени Н.Э.Баумана \\*
        \vspace{-12mm}
        \begin{figure}[h]
            \center{\includegraphics[width=0.2\linewidth]{symbol}}
        \end{figure}
        \vspace{-16mm}
        \hrulefill
    \end{center}
    \center{Факультет «Робототехника и комплексная автоматизация»\\
    		Кафедра «Системы Автоматизированного Проектирования»}
    \begin{center}
        \Large Отчёт о преддипломной практике по теме:
    \end{center}
    
    \vspace{2.5em}
    
    \begin{center}
        \textsc{\textbf{Инструментарий коллективной оптимизации программ}}
    \end{center}
    
    \vspace{6em}
    
    \begin{flushleft}
        \hspace{8.5cm}Исполнитель \hrulefill Панков М.К. \\
        \vspace{1.5em}
        \hspace{8.5cm}Руководитель \hrulefill Карпенко А.П.\\
    \end{flushleft}
    
    \vspace{\fill}
    
    \begin{center}
        Москва, \\*
        2013
    \end{center}

\end{titlepage}
% Учебное задание
\setcounter{page}{2}
\section*{Учебное задание}
% TODO
Разработать инструментарий анализа производительности программ. Инструментарий должен включать в себя компоненты, перечисленные ниже.
\begin{itemize}
    \item Репозиторий исходных кодов программ.
    \item База данных экспериментов.
    \item Интерактивное текстовое окружение для проведения экспериментов, обработки и сохранения данных.
    \item Компонент построения графиков для визуального представления данных экспериментов.
\end{itemize}

\pagebreak
% Содержание
\tableofcontents
\pagebreak
% Введение
\chapter*{Введение}
\addcontentsline{toc}{chapter}{Введение}%
Разработка современного математического обеспечения САПР "--- сложная задача, требующая больших вложений материальных и временных ресурсов. Зачастую для достижения оптимальности по таким критериям, как производительность, объём занимаемой памяти, и энергопотребление требуется ручная настройка компилятора "--- выбор вариантов тонких настроек оптимизации под определённую программно-аппаратную платформу с учётом особенностей решаемой задачи. Это плохо формализованная задача, которая часто решается разработчиками МО методом проб и ошибок. В результате этого оптимальность часто не достигается ни по одному из выбранных критериев.

Распространённой проблемой для исследователей и инженеров является то, что параметры компиляции подбираются под конкретный набор входных данных, на конкретной программно-аппаратной платформе \cite{Jimenez:2008:PRC:1505816.1505822,luo:inria-00436034}. При этом для набора данных, встречающегося в реальных задачах (а не используемого во время подбора оптимальных параметров), оптимальные параметры могут оказаться иными.

Для увеличения возможностей по тонкой оптимизации математического программного обеспечения для рабочих станций и суперкомпьютеров необходима формализация поисковой области. С помощью построения модели производительности программ возможно достичь лучшего понимания воздействия оптимизаций компилятора на интересующие разработчика критерии эффективности программы. Таким образом можно сделать поиск оптимальных настроек более направленным и локализованным, сокращая цикл разработки математического обеспечения, стоимость разработки и поддержки.

Цель — разработать систему сбора, систематизации, формализации данных о производительности компилируемых программ в зависимости от настроек компилятора и программно-аппаратной платформы, а также выполняющую функции поддержки базы знаний и обучения с целью моделирования и предсказания эффективности программ.

Примером критерия эффективности является производительность программы на данной программно-аппаратной платформе.

\chapter{Конструкторская часть}
\section{Статистический анализ быстродействия программ}
\subsection{Технологии компиляторов}
Современный оптимизирующий промышленный компилятор "--- сложная программа, состоящая из нескольких модулей, типично выделяемых в две части~\cite{Aho:2006:CPT:1177220}, описанные ниже.
\begin{enumerate}
	\item Анализирующая часть.
	\begin{enumerate}
		\item Лексический анализатор. Распознаёт строки (<<лексемы>>) и формирует из них типизированные значения (<<токены>>).
		\item Синтаксический анализатор. Потребляет токены и анализирует грамматическую правильность сформированных из них конструкций.
		\item Семантический анализатор. Производит семантический анализ --- оценивает правильность и осмысленность сформированных из токенов конструкций.
	\end{enumerate}
	\item Синтезирующая часть.
	\begin{enumerate}
		\item Генератор абстрактного синтаксического дерева. Производит построение абстрактного синтаксического дерева из токенов. Дерево отражает связи между различными элементами исходного файла, но уже не привязано к их строковому представлению. Оно содержит множество типизированных узлов со связями между ними.
		\item Генератор промежуточного представления. Создаёт промежуточное представление (такое как трёхадресный код или форма с статическим единственным присваиванием) из абстрактного синтаксического дерева. Также может конвертировать одно промежуточное представление в другое для удобства произведения оптимизаций.
		\item Оптимизатор. Производит трансформацию промежуточного представления с сохранением семантики и увеличением эффективности. Типичными оптимизациями являются уменьшение сложности операции и разворачивание циклов. Обычно использует набор эвристик, определяющих правила применения оптимизаций и их параметры, т.к. многие оптимизации имеют настройку того или иного рода, определяющие её агрессивность.
		\item Генератор исполняемого файла. Производит окончательное создание исполняемого файла в двоичном формате: перевод промежуточного представления в набор инструкций целевой машины, оптимизации на объектном файле и связывание объектных файлов в исполняемый.
	\end{enumerate}
\end{enumerate}

Стоит отметить, что данное деление является довольно условным: так, синтаксический анализатор, семантический анализатор и генератор абстрактного синтаксического дерева типично объединяются в одну единицу, называемую <<парсер>>. Однако, оно довольно показательно с точки зрения объёма и разнообразия задач, решаемых компилятором.


\subsection{Сложность разработки компиляторов}
Разработка современного оптимизирующего промышленного компилятора "--- сложная задача, требующая огромных вложений человеческих ресурсов. Больш\'{а}я часть этой сложности обусловлена необходимостью тонкой настройки эвристических методов оптимизации исполняемого кода (которые в общем случае решаются NP-сложными алгоритмами), производимой вручную путём просмотра генерируемого кода и состояний промежуточного представления компилятора. Не говоря о том, что само по себе создание таких эвристик представляет собой зачастую нетривиальную задачу \cite{Muchnick:2004:EIS:989393.989413,Lavery:1995:IPC:626512.626987,Abraham:1996:MSR:243846.243903}, эвристики часто имеют противоположные цели "--- например, одни оптимизации увеличивают производительность кода ценой размера исполняемого файла (например, разворачивание циклов), а другие уменьшают размер кода, уменьшая производительность. Помимо этого, в компиляторах чаще всего используются упрощённые модели аппаратного обеспечения, не учитывающие многие эффекты даже первого порядка, не говоря уже о побочных эффектах. Более того, сами правила обнаружения случаев, в которых можно оптимизировать код, создаются человеком, что также трудоёмко и является далеко не исчерпывающим способом увеличения производительности результирующего кода. Помимо этого, от порядка произведения оптимизаций часто зависит итоговая производительность и размер кода исполняемого файла.

В то же время, у настроек эвристик есть фокальные точки "--- такие настройки оптимизации, при которых достигается оптимум по какому-либо критерию. Эти фокальные точки могут быть найдены с помощью функции цели. Примером задачи, для которой существуют достаточно хорошие целевые функции, является задача распределения регистров.

В существующих компиляторах также стоит проблема подбора оптимальных настроек компиляции для данной конкретной программы. Чаще всего проще ограничиться стандартной настройкой <<максимальная производительность>>, нежели производить трудоёмкий анализ кода, который так же требует квалификации опытного разработчика компилятора и знания о платформе, для которой генерируется исполняемый файл. Эта задача также может решаться автоматически и, в ряде случаев, более эффективно.


\subsection{Подходы к решению проблемы сложности эвристик}
Попытки настраивать и даже синтезировать \cite{Stephenson:2006:ACC:1269626,Stephenson:2003:MOI:780822.781141} эвристики предпринимались неоднократно. Как уже показали работы \cite{Agakov:2006:UML:1121992.1122412,Bodin98iterativecompilation,FCA2007,Cooper:2005:AAC:1065910.1065921}, задача подбора самих эвристик и их настройки может решаться автоматически.

Принципиально, классическим подходом к этой задаче является итеративная компиляция. Также, в последнее время получили распространение методы машинного обучения и коллективной оптимизации. Далее мы рассмотрим эти подходы.


\subsection{Итеративная компиляция}
Итеративная компиляция "--- методика поиска оптимальных по некоторому критерию настроек компилятора, при которой компилируемая программа собирается много раз с одним набором данных на одной и той же программно-аппаратной платформе с разными, выбираемыми, как правило, случайным образом, настройками, с последующим запуском и замером интересующих факторов во время её выполнения (например, времени выполнения). Очевидно, что такой подход страдает от чудовищной неэффективности --- поиск оптимальных настроек для одной программы может занимать месяц машинного времени. Количество необходимых компиляций составляет десятки и сотни запусков \cite{Kisuki:1999:FSI:646347.690219}. Кроме того, такой поиск является весьма ограниченным, поскольку найденные настройки в итоге подходят только для данного набора данных на данной платформе с данным компилятором (вплоть до конкретной версии), и практически никаких общих рекомендаций извлечь из этого процесса не получается.

Неоднократно производились попытки ускорить этот процесс с помощью применения различных методов оптимизации, изменяющих характер поиска настроек с случайного на более направленный. В этом случае задача осложняется чрезвычайно большой многомерностью такой оптимизации "--- признаков, по которым опознаются программы и которые можно считать аргументами функции эффективности, насчитываются десятки (имеются ввиду такие признаки, собираемые с исходного кода программы, как средняя глубина вложенности циклов, среднее число вызовов процедур на блок и т.д.). Таким образом, данный способ решения является крайне плохо приспособленным для реального применения.


\subsection{Машинное обучение}
Недавние работы в данной области показали возможность увеличить эффективность итеративной компиляции с помощью машинного обучения, позволяя обеспечить компилятор возможностью подстраиваться под сложные архитектуры и повысить его производительность относительно классических статических компиляторов "--- т.е. таких, которые не были специально модифицированы для работы с инструментарием машинного обучения \cite{Dubach:2009:PCO:1669112.1669124,Dubach:2008:EPA:1450095.1450103}.

Далее для понимания связи машинного обучения с компиляцией приведена классическая постановка задачи этого раздела науки.

Имеется множество объектов (ситуаций) и множество возможных ответов (откликов, реакций). Существует некоторая зависимость между ответами и объектами, но она неизвестна. Известна только конечная совокупность прецедентов "--- пар <<объект, ответ>>, называемая обучающей выборкой. На основе этих данных требуется восстановить зависимость, то есть построить алгоритм, способный для любого объекта выдать достаточно точный ответ. Для измерения точности ответов определённым образом вводится функционал качества.

В нашем случае объектом является исходный код программы с его признаками (такими, как средний размер массивов и среднее число итераций циклов), а ответом "--- значение интересующего нас фактора во время исполнения программы (например, объём занимаемой памяти).

Стоит отметить, что предпочтительным является обучение с использованием только производственных запусков компилятора, т.е. без специальных обучающих запусков. Для этого предлагается использовать коллективную оптимизацию, позволяющую пользователям обмениваться данными об успешной оптимизации и не проделывать избыточную работу каждый раз изолированно.


\subsection{Коллективная оптимизация}
Коллективная оптимизация предполагает реализацию некого инструментария, подключаемого к компилятору с целью организации его прозрачной работы с общим репозиторием настроек компиляции. Этот подход является многообещающим, но не лишён сложностей.

А именно, перед реализацией коллективной оптимизации стоит несколько проблем исследовательского и инженерного характера. Инженерная проблема состоит в создании системы, позволяющей прозрачно переиспользовать накопленные другими пользователями знания о произведенных компиляциях и отправлять данные о них без дополнительных действий со стороны конечного пользователя компилятора. Частью решения этой проблемы является плагин к компилятору GCC, предоставляющий возможности произвольного включения выбранных проходов компилятора и обеспечивающий его взаимодействие с репозиторием кода cTuning.org \cite{Fur2009}. В случае более современного компилятора llvm \cite{llvm} плагин имеет более простую организацию, поскольку данный инструмент имеет встроенные возможности произвольного выбора набора производимых оптимизаций, и нужно обеспечить только взаимодействие с репозиторием cTuning. Исследовательская проблема состоит в изучении воздействия различных оптимизаций на различные факторы получаемых программ, такие, как время исполнения и объём используемой памяти.

Коллективная оптимизация может обеспечить эффективность итеративной компиляции без произведения такого большого количества избыточных запусков компилятора. Одним из успешно применяемых подходов к этой задаче предполагает статистическое сравнение эффективности пар наборов оптимизаций. Недавние работы также обнаружили значительно б\'{о}льшую важность обучения на различных наборах данных и на различных архитектурах, нежели обучения на различных программах \cite{springerlink:10.1007/11596110:24}. Это опровергает доминировавшую до этого точку зрения на применение машинного обучения в компиляторах и предлагает новые перспективы в области автоматической настройки компиляции.


\subsection{Выводы}
Существующие методы статистического анализа быстродействия программ плохо применимы на практике ввиду отсутствия удобного инструментария для проведения исследований и большого времени, требуемого для проведения запусков программ.

Наиболее перспективными подходами к решению задачи анализа быстродействия являются машинное обучение предсказателей производительности и коллективное принятие решений в сфере анализа и оптимизации быстродействия.
% Обзор области
\chapter{Обзор области автоматической настройки компилятора}
\section{Технологии компиляторов}
Современный оптимизирующий промышленный компилятор --- сложная программа, состоящая из нескольких модулей, типично выделяемых в 2 части \cite{Aho:2006:CPT:1177220}:
\begin{enumerate}
	\item Анализирующая часть
	\begin{enumerate}
		\item Лексический анализатор. Распознаёт строки («лексемы») и формирует из них типизированные значения («токены»).
		\item Синтаксический анализатор. Потребляет токены и анализирует грамматическую правильность сформированных из них конструкций.
		\item Семантический анализатор. Производит семантический анализ --- оценивает правильность и осмысленность сформированных из токенов конструкций.
	\end{enumerate}
	\item Синтезирующая часть
	\begin{enumerate}
		\item Генератор абстрактного синтаксического дерева. Производит построение абстрактного синтаксического дерева из токенов. Дерево отражает связи между различными элементами исходного файла, но уже не привязано к их строковому представлению. Оно содержит множество типизированных узлов со связями между ними.
		\item Генератор промежуточного представления. Создаёт промежуточное представление (такое как трёхадресный код или форма с статическим единственным присваиванием) из абстрактного синтаксического дерева. Также может конвертировать одно промежуточное представление в другое для удобства произведения оптимизаций.
		\item Оптимизатор. Производит трансформацию промежуточного представления с сохранением семантики и увеличением эффективности. Типичными оптимизациями являются уменьшение сложности операции и разворачивание циклов. Обычно использует набор эвристик, определяющих правила применения оптимизаций и их параметры, т.к. многие оптимизации имеют настройку того или иного рода, определяющие её агрессивность.
		\item Генератор исполняемого файла. Производит окончательное создание исполняемого файла в двоичном формате: перевод промежуточного представления в набор инструкций целевой машины, оптимизации на объектном файле и связывание объектных файлов в исполняемый.
	\end{enumerate}
\end{enumerate}

Стоит отметить, что данное деление является довольно условным: так, синтаксический анализатор, семантический анализатор и генератор абстрактного синтаксического дерева типично объединяются в одну единицу, называемую «парсер». Однако, оно довольно показательно с точки зрения объёма и разнообразия задач, решаемых компилятором.
\section{Сложность разработки компиляторов}
Разработка современного оптимизирующего промышленного компилятора --- сложная задача, требующая огромных вложений человеческих ресурсов. Больш\'{а}я часть этой сложности обусловлена необходимостью тонкой настройки эвристических методов оптимизации исполняемого кода (которые в общем случае решаются NP-сложными алгоритмами), производимой вручную путём просмотра генерируемого кода и состояний промежуточного представления компилятора. Не говоря о том, что само по себе создание таких эвристик представляет собой зачастую нетривиальную задачу \cite{Muchnick:2004:EIS:989393.989413,Lavery:1995:IPC:626512.626987,Abraham:1996:MSR:243846.243903}, эвристики часто имеют противоположные цели --- например, одни оптимизации увеличивают производительность кода ценой размера исполняемого файла (например, разворачивание циклов), а другие уменьшают размер кода, уменьшая производительность. Помимо этого, в компиляторах чаще всего используются упрощённые модели аппаратного обеспечения, не учитывающие многие эффекты даже первого порядка, не говоря уже о побочных эффектах. Более того, сами правила обнаружения случаев, в которых можно оптимизировать код, создаются человеком, что также трудоёмко и является далеко не исчерпывающим способом увеличения производительности результирующего кода. Помимо этого, от порядка произведения оптимизаций часто зависит итоговая производительность и размер кода исполняемого файла.

В то же время, у настроек эвристик есть фокальные точки --- такие настройки оптимизации, при которых достигается оптимум по какому-либо критерию. Эти фокальные точки могут быть найдены с помощью функции цели. Примером задачи, для которой существуют достаточно хорошие целевые функции, является задача распределения регистров.

В существующих компиляторах также стоит проблема подбора оптимальных настроек компиляции для данной конкретной программы. Чаще всего проще ограничиться стандартной настройкой «максимальная производительность», нежели производить трудоёмкий анализ кода, который так же требует квалификации опытного разработчика компилятора и знания о платформе, для которой генерируется исполняемый файл. Эта задача также может решаться автоматически и, в ряде случаев, более эффективно.

\section{Подходы к решению проблемы сложности эвристик}
Попытки настраивать и даже синтезировать \cite{Stephenson:2006:ACC:1269626,Stephenson:2003:MOI:780822.781141} эвристики предпринимались неоднократно. Как уже показали работы \cite{Agakov:2006:UML:1121992.1122412,Bodin98iterativecompilation,FCA2007,Cooper:2005:AAC:1065910.1065921}, задача подбора самих эвристик и их настройки может решаться автоматически.

Принципиально, классическим подходом к этой задаче является итеративная компиляция. Также, в последнее время получили распространение методы машинного обучения и коллективной оптимизации. Далее мы рассмотрим эти подходы.

\section{Итеративная компиляция}
Итеративная компиляция --- методика поиска оптимальных по некоторому критерию настроек компилятора, при которой компилируемая программа собирается много раз с одним набором данных на одной и той же программно-аппаратной платформе с разными, выбираемыми, как правило, случайным образом, настройками, с последующим запуском и замером интересующих факторов во время её выполнения (например, времени выполнения). Очевидно, что такой подход страдает от чудовищной неэффективности --- поиск оптимальных настроек для одной программы может занимать месяц машинного времени. Количество необходимых компиляций составляет десятки и сотни запусков \cite{Kisuki:1999:FSI:646347.690219}. Кроме того, такой поиск является весьма ограниченным, поскольку найденные настройки в итоге подходят только для данного набора данных на данной платформе с данным компилятором (вплоть до конкретной версии), и практически никаких общих рекомендаций извлечь из этого процесса не получается.

Неоднократно производились попытки ускорить этот процесс с помощью применения различных методов оптимизации, изменяющих характер поиска настроек с случайного на более направленный. В этом случае задача осложняется чрезвычайно большой многомерностью такой оптимизации --- признаков, по которым опознаются программы и которые можно считать аргументами функции эффективности, насчитываются десятки (имеются ввиду такие признаки, собираемые с исходного кода программы, как средняя глубина вложенности циклов, среднее число вызовов процедур на блок и т.д.). Таким образом, данный способ решения является крайне плохо приспособленным для реального применения.

\section{Машинное обучение}
Недавние работы в данной области показали возможность увеличить эффективность итеративной компиляции с помощью машинного обучения, позволяя обеспечить компилятор возможностью подстраиваться под сложные архитектуры и повысить его производительность относительно классических статических компиляторов --- т.е. таких, которые не были специально модифицированы для работы с инструментарием машинного обучения \cite{Dubach:2009:PCO:1669112.1669124,Dubach:2008:EPA:1450095.1450103}.

Далее для понимания связи машинного обучения с компиляцией приведена классическая постановка задачи этого раздела науки.

Имеется множество объектов (ситуаций) и множество возможных ответов (откликов, реакций). Существует некоторая зависимость между ответами и объектами, но она неизвестна. Известна только конечная совокупность прецедентов — пар «объект, ответ», называемая обучающей выборкой. На основе этих данных требуется восстановить зависимость, то есть построить алгоритм, способный для любого объекта выдать достаточно точный ответ. Для измерения точности ответов определённым образом вводится функционал качества.

В нашем случае объектом является исходный код программы с его признаками (такими, как средний размер массивов и среднее число итераций циклов), а ответом --- значение интересующего нас фактора во время исполнения программы (например, объём занимаемой памяти).

Стоит отметить, что предпочтительным является обучение с использованием только производственных запусков компилятора, т.е. без специальных обучающих запусков. Для этого предлагается использовать коллективную оптимизацию, позволяющую пользователям обмениваться данными об успешной оптимизации и не проделывать избыточную работу каждый раз изолированно.

\section{Коллективная оптимизация}
Коллективная оптимизация предполагает реализацию некого инструментария, подключаемого к компилятору с целью организации его прозрачной работы с общим репозиторием настроек компиляции. Этот подход является многообещающим, но не лишён сложностей.

А именно, перед реализацией коллективной оптимизации стоит несколько проблем исследовательского и инженерного характера. Инженерная проблема состоит в создании системы, позволяющей прозрачно переиспользовать накопленные другими пользователями знания о произведенных компиляциях и отправлять данные о них без дополнительных действий со стороны конечного пользователя компилятора. Частью решения этой проблемы является плагин к компилятору GCC, предоставляющий возможности произвольного включения выбранных проходов компилятора и обеспечивающий его взаимодействие с репозиторием кода cTuning.org \cite{Fur2009}. В случае более современного компилятора llvm \cite{llvm} плагин имеет более простую организацию, поскольку данный инструмент имеет встроенные возможности произвольного выбора набора производимых оптимизаций, и нужно обеспечить только взаимодействие с репозиторием cTuning. Исследовательская проблема состоит в изучении воздействия различных оптимизаций на различные факторы получаемых программ, такие, как время исполнения и объём используемой памяти.

Коллективная оптимизация может обеспечить эффективность итеративной компиляции без произведения такого большого количества избыточных запусков компилятора. Одним из успешно применяемых подходов к этой задаче предполагает статистическое сравнение эффективности пар наборов оптимизаций. Недавние работы также обнаружили значительно б\'{о}льшую важность обучения на различных наборах данных и на различных архитектурах, нежели обучения на различных программах \cite{springerlink:10.1007/11596110:24}. Это опровергает доминировавшую до этого точку зрения на применение машинного обучения в компиляторах и предлагает новые перспективы в области автоматической настройки компиляции.
% Основная часть
% Мотивация, задачи, план
\chapter{Основная часть}
\section{Мотивация}
\begin{enumerate}
\item Многие исследования показывают потенциал, но плохо систематизированы.
\item Инженерная задача по созданию инфраструктуры для коллективной оптимизации не решена до сих пор.
\item Конечным пользователям компиляторов мешает именно отсутствие инструментария, который помог бы им пользоваться обучаемым компилятором так же, как обычным (интерфейс чёрного ящика).
\item В целом это популярная в последнее время тема — Intel запустил Exascale Lab, также IBM и ARM очень заинтересованы и хотят финансировать эти исследования.
\end{enumerate}

\section{Цель}
Цель — разработать систему сбора, систематизации, формализации данных о производительности компилируемых программ в зависимости от настроек компилятора, а также выполняющую функции поддержки базы знаний и обучения с целью:
\begin{enumerate}
\item Предсказания производительности программ
\item Оптимизации производительности программ путём изменения настроек компилятора
\end{enumerate}

Пока рассматриваем только 1 критерий --- производительность.

\subsection{Задачи}
\begin{enumerate}
\item {Выполнить первоначальное наполнение репозитория данными.
Установить окружение и репозиторий локально и внести туда программы из центрального репозитория.
Собрать небольшое количество признаков для начала работы.\\
Предполагаемый объём работы от общего: 5\%.}
\item {Пронаблюдать влияние набора данных, компилятора и архитектуры на производительность. 
Для этого попробовать перебор разных комбинаций этих параметров (в числе которых и настройки компилятора) и сравнить получаемые результаты. Это нужно для набора данных для последующего машинного обучения на них. 
Это не полный перебор, а небольшой просмотр области поиска с целью выяснения реакции программ на изменения для последующей формализации модели производительности.\\
Предполагаемый объём работы от общего: 10\%.}
\item {Формализовать признаки программ, влияние платформы и настроек компилятора на производительность.
Строим модель производительности программы. Это будет делаться изначально вручную с последующими попытками обобщить до автоматического выбора модели с помощью машинного обучения. Пример: начать с линейной регрессии от одного признака (размера входных данных), дальше обнаруживать нестыковки и добавлять новые признаки.
Полезными моделями могут быть вероятностная (она удобна для представления сравнения 2 программ и легко интерпретируется для понимания человеком), деревья решений (легко интерпретируется, относительно быстро строится и работает).
Другие методы, такие как SVM и kNN, могут тоже быть полезны и возможно будут использоваться — но они лишены такой черты, как простота понимания, почему именно такое решение было выбрано моделью. Это затрудняет дальнейшее развитие проекта (поскольку непонятно, в каком месте модели произошла ошибка). Также, применение методов “ad-hoc” не даёт воспроизводимых результатов, на которых можно хорошо строить модель.\\
Предполагаемый объём работы от общего: 35\%.}
\item {Построить предсказывающий модуль
Он может работать в 2 режимах:
	\begin{enumerate}
	\item Предсказание уровня производительности программы при заданных входных параметрах. Полезно в облачных сервисах и параллельных вычислениях (балансировка загрузки). Это регрессионный анализ.
	\item Выбор оптимальных настроек компилятора для получения максимальной производительности при заданной платформе, программе (которой ещё нет в репозитории известных программ) и множестве наборов входных данных. Возможны даже рекомендации по выбору другого компилятора. Это классификация.
	\end{enumerate}
Сюда же попадает первоначальное обучение предсказателя для начала работы системы.\\
Предполагаемый объём работы от общего: 25\%.}
\item {Анализировать проблемы предсказателя (ошибки неправильной классификации или регрессии) с целью улучшить работу.
Это включает в себя как ручной просмотр моделей экспертами с помощью нашего интрументария (без него это делать крайне неудобно) с целью понимания возможной проблемы, так и попытки автоматически улучшить модель с помощью добавления новых признаков или настроек, не учтённых в модели для этого.
Сюда же попадает дообучение при получении новых данных, таких как запуск программы с новыми данными, на новой платформе или с новым компилятором, или данные о новых программах.\\
Предполагаемый объём работы от общего: 20\%}
\item {Улучшить интеграцию в существующие компиляторы, так чтобы для конечного пользователя работа системы происходила прозрачно (интерфейс работы с компилятором не изменялся). Это нужно для вовлечения большего числа разработчиков и экспертов в компьютерных системах для выявления проблем моделей на этапе 5.\\
Предполагаемый объём работы от общего: 5\%.}
\end{enumerate}

Реализация будет использовать Python в качестве основного языка с возможными вызовами scipy и R для реализации сложной обработки данных. Это простой и удобный, популярный и хорошо поддерживаемый язык с большим количеством библиотек. В случае необходимости можно достигнуть высокой производительности с помощью компилятора PyPy.
\section{Введение}
\subsection{Мотивация}
Ниже приведены основные причины заинтересованности в данной задаче.
\begin{itemize}
    \item Многие исследования в области применения машинного обучения в компиляции \cite{Dubach2009PCO,Thomson2009RTT,Stephenson2003MOI} показывают большой потенциал и востребованность решения данной задачи, но плохо систематизированы.
    \item Инженерная задача по созданию инфраструктуры для коллективной оптимизации не решена до сих пор \cite{Fursin2010COP}.
    \item Конечным пользователям компиляторов мешает именно отсутствие инструментария, который помог бы им пользоваться обучаемым компилятором так же, как обычным (используя интерфейс «чёрного ящика»).
    \item Производители аппаратного и программного обеспечения, такие как Intel, IBM, ARM заинтересованы в решении такой задачи \cite{IntelExascale}.
\end{itemize}


\subsection{Цель}
Конечная цель, которой требуется достичь --- разработать систему сбора, систематизации, формализации данных о производительности компилируемых программ в зависимости от настроек компилятора, а также выполняющую функции поддержки базы знаний и обучения. Цели данной системы представлены ниже.
\begin{itemize}
    \item Предсказание производительности программ.
    \item Оптимизация производительности программ путём изменения настроек компилятора.
\end{itemize}

На данном, начальном, этапе работы рассматривается только один критерий оптимальности --- производительность, а именно --- полное время выполнения программы.

Однако, сравнивая стоящую перед нами задачу с текущими успехами в этой области исследований \cite{Fursin2010COP}, обнаруживаем, что достижение данной цели в рамках курсового проекта не представляется возможным. В связи с этим разработка ограничена учебным заданием и продолжение работы будет осуществляться в рамках дипломного проекта.

Итак, цель данной работы --- разработать основные компоненты расширяемого инструментария для проведения экспериментов по оптимизации программ (в ручном и автоматическом режиме) и прозрачного использования данных из экспериментов в промышленной компиляции.


\subsection{Обзор предшествующих работ}
Как было сказано выше, многие исследования в области применения машинного обучения в компиляции \cite{Dubach2009PCO,Thomson2009RTT,Stephenson2003MOI} показывают большой потенциал и востребованность решения данной задачи, но плохо систематизированы. Вследствие этого основной проблемой является реализация общей базы данных программ и экспериментов и каркаса инструментария для проведения исследований. Реализация всех конечных сценариев, необходимых в исследовании, не является приоритетной задачей, но должно быть осуществимо достаточно простым способом силами самих программных инженеров-конечных пользователей системы и других исследователей.

Ниже приведены некоторые отличия данной работы от существующих.
\begin{itemize}
    \item Исследователям и инженерам нужен в первую очередь минимальный функционал для решения описываемой задачи. Графический интерфейс не является необходимым. В сосредоточенности на традиционном текстовом командном интерфейсе основное отличие этой работы от работы \cite{Fursin2010COP}.
    \item Инструментарий не может реализовывать все необходимые сценарии работы с самого начала. Поэтому цель --- сделать систему компактной, модульной и документированной, что обеспечит хорошую поддерживаемость и расширяемость. В этом также одно из отличий от работы \cite{Fursin2010COP}.
    \item Разработка системы сфокусирована на применении локально в изоляции от других исследователей, поскольку это --- первоначально необходимый сценарий использования. Разработка многопользовательских сценариев является следующим шагом. Однако используемая технологическая база позволяет относительно просто расширить инструментарий для использования в распределённых системах.
\end{itemize}


\subsection{Задачи}
Задачи, решаемые в данной работе, перечислены далее в примерном порядке реализации.
\begin{enumerate}
    \item Установить репозиторий исходных кодов программ.
    \item Создать и настроить базу данных для хранения экспериментов.
    \item Реализовать интерактивное текстовое окружение для проведения экспериментов, обработки и сохранения данных.
    \item Реализовать построение графиков для визуального представления данных экспериментов.
\end{enumerate}
\pagebreak
% Теоретическая часть
\part{Теоретическая часть}

\section{Регрессионный анализ}

Hello \cite{lag}.
% Исследовательская часть
\section{Исследовательская часть}

% Инструментарий коллективной оптимизации
\section{Конструкторская часть}

За основу взята система, разработанная в рамках курсового проекта \cite{adaptor}. Представление эксперимента доработано так, как описано в следующем разделе.

\subsection{Структура набора тестовых программ Polybench}

Структура каталогов с исходными кодами тестовых программ представляет собой дерево. Типичный путь до папки исходных кодов выглядит так: 'polybench/linear-algebra/kernels/symm'.

Для обеспечения простоты использования инструментария был реализован поиск исходных кодов заданной программы в дереве каталогов. При начале работы с системой пользователь лишь задаёт название программы, которую он хочет исследовать, а инструментарий находит исходные коды автоматически.

Для этого был реализован поиск по дереву путей путём обхода в ширину. Поиск останавливается, когда имя каталога, имеющего наибольший уровень вложенности, совпадает с заданным пользователем названием программы.

\subsection{Эксперимент по оценке производительности программы}

Эксперимент по оценке производительности программы представляет собой последовательность шагов, перечисленную ниже.

\begin{enumerate}
	\item Сборка программы.
	\item Измерение времени исполнения программы.
	\item Сохранение данных об эксперименте в базу данных.
\end{enumerate}

Далее, если не указано иное, мы будем называть эксперимент по оценке производительности программы просто «экспериментом».

\subsection{Реализация модели производительности программы}
\label{perf_mod_impl}

Модель производительности программы строится по множеству экспериментов. В системе эксперимент характеризуется свойствами, приведёнными ниже. Многие свойства не являются непосредственно ценными для анализа производительности на данный момент, но сохранены в структуре эксперимента с целью полноты описания и возможности воспроизведения эксперимента другим исследователем --- потенциально на другом компьютере.

\begin{itemize}
	\item Название программы.
	\item Настройки сборки. Включают в себя следующие свойства:
	\begin{itemize}
	    \item используемый компилятор. Задаётся в виде команды запуска компилятора из командной строки;
	    \item базовый уровень оптимизации. Задаётся в виде строки типа «-O<число>», где <число> означает уровень оптимизации. Допустимы значения от 1 до 3;
	    \item дополнительные настройки оптимизации. Задаются в виде строки «-f<имя настройки>», например «-floop-unroll»;
	    \item дополнительные настройки компиляции. Задаются в виде строки, содержащей настройки путей поиска заголовочных файлов («-I<путь>») и библиотек(«-L<путь>»);
	    \item дополнительные настройки связывания. Задаются в виде строки, содержащей список необходимых библиотек («-l<имя библиотеки>»);
	    \item каталог с исходными текстами тестовых программ;
	    \item путь к главному исходному файлу тестовой программы.
	\end{itemize}

	\item Результат запуска с калибровкой. Включает в себя следующие свойства:
	\begin{itemize}
		\item суммарное время исполнения программы;
		\item число запусков программы;
		\item вычисленное время однократного исполнения программы;
		\item дисперсия суммарного времени исполнения;
		\item вариация суммарного времени исполнения;
		\item список суммарных времён исполнения.
	\end{itemize}

	\item Результат проверки. Включает в себя следующие свойства:
	\begin{itemize}
		\item измеренное время однократного исполнения программы;
		\item действительное время однократного исполнения программы (если заведомо известно);
		\item абсолютная ошибка измерения (если известно действительное время однократного исполнения);
		\item относительная ошибка измерения (если известно действительное время однократного исполнения).
	\end{itemize}

	\item Информация об аппаратном обеспечении, на котором исполняется экспериментальная программа. Включает в себя следующую информацию о процессоре:
	\begin{itemize}
		\item текстовое название процессора;
		\item частота процессора;
		\item размер кэша процессора;
		\item список возможностей процессора.
	\end{itemize}

	\item Информация о вычислительной задаче. На данный момент включает в себя два свойства:
	\begin{itemize}
		\item число столбцов в обрабатываемой матрице;
		\item число строк в обрабатываемой матрице.
	\end{itemize}
\end{itemize}

\subsection{Анализ производительности программы}

Ниже приведены непосредственно важные для анализа производительности свойства эксперимента. Это множество свойств является подмножеством всех свойств эксперимента, описанных в разделе \ref{perf_mod_impl}.

\begin{itemize}
	\item Название программы.
	\item Используемый компилятор.
	\item Используемый базовый уровень оптимизации.
	\item Используемые дополнительные настройки оптимизации.
	\item Число строк в обрабатываемой матрице.
	\item Число столбцов в обрабатываемой матрице.
	\item Измеренное время однократного исполнения программы.
	\item Название процессора, на котором исполнялась программа.
	\item Частота процессора.
	\item Размер кэша процессора.
	\item Множество дополнительных возможностей, поддерживаемых или не поддерживаемых процессором. Примером такого свойства является наличие поддержки процессором набора инструкций SSE \cite{sse}.
\end{itemize}

\subsection{Централизованное хранилище экспериментов}

\subsubsection{Модуль работы с базой данных}

Для обмена данными с централизованным хранилищем экспериментов реализован модуль работы с базой данных CouchDB. Он основан на библиотеке couchdbkit \cite{couchdbkit} для языка программирования Python и реализует документы (классы сущностей, хранимых в БД), описанные ниже. Документы хранят свойства соответствующих сущностей в формате, пригодном для хранения в БД. Состав документа индентичен составу соотвествующей сущности, как описано в предыдущем разделе.

\begin{itemize}

\item Документ эксперимента.
\item Документ результата измерения времени исполнения программы.
\item Документ результата проверки точности измерения времени исполнения программы.
\item Документ информации об аппаратном обеспечении.
\item Документ настроек сборки программы.

\end{itemize}

\subsubsection{Размещение базы данных на удалённом сервере}

В качестве удалённого сервера выбран сервер Amazon Elastic Computing (иначе называемый Amazon EC2) \cite{amazon-ec2}. Amazon EC2 предоставляет услуги размещения ресурсов в интернет по схеме так называемых «облачных вычислений» (\textit{англ. cloud computing} \cite{cloud-computing}). Получение доступа к серверу, установка на него операционной системы, включение и выключение, подсчёт стоимости автоматизированы и выполняются через веб-интерфейс. Стоимость предоставления услуг является очень низкой при текущих объёмах использования вычислительных мощностей, хранилища и передаваемых через сеть данных \cite{amazon-billing}. Объём предоставляемых услуг легко поддаётся масштабированию при увеличении нагрузки.

Для обеспечения возможностей централизованного сбора данных и их анализа в фоновом режиме (\textit{англ. offline}) был реализован серверный модуль для работы на платформе Amazon EC2. Данный модуль может запускать сценарии анализа накопленных экспериментов по расписанию или при наступлении каких-либо других условий --- например, по достижении определённого количества сохранённых в БД экспериментов. Это необходимо для постоянного улучшения модели производительности программ на основе поступающих данных.

Серверный модуль также может выполнять эксперименты непосредственно на аппаратном обеспечении сервера. На данный момент это полезно и допустимо, поскольку нагрузка на сервер, связанная с обработкой клиентских запросов, пренебрежимо мала, а данные о запусках программ на ещё одной аппаратной платформе имеют большую ценность в связи с ограниченными материальными ресурсами.

Данный модуль также поддерживает работу как с удалённой БД на платформе CloudAnt, так и с локальной (относительно сервера) базой данных. А именно, база данных CouchDB может быть установлена на том же физическом сервере Amazon EC2, что обеспечивает более низкие задержки передачи данных между серверным модулем и БД. Переключение испоьзуемой БД производится изменением одного параметра настройки системы.

Отдельно следует отметить высокую совместимость клиентского и серверного модулей инструментария. Оба они могут работать как с локальной, так и с удалённой БД. Оба модуля поддерживают одно и то же множество сценариев и предоставляют пользователю одинаковые команды. Возможна двусторонняя синхронизация данных между базами данных CouchDB, расположенными на разных физических машинах.

Перспективной является возможность реализации API для работы с системой как с веб-сервисом. На данный момент API не реализовано ввиду отсутствия необходимости. Наличие API имеет смысл при большом спросе на доступ к системе со стороны клиентов.

Сочетание использования сервера Amazon EC2 и базы данных CouchDB делает систему легко масштабируемой и позволяет легко построить веб-сервис по модели «облачных вычислений».

\subsection{Сценарии исследования}

Все сценарии исследования подразумевают многократный запуск одной и той же программы на одном компьютере с разными размерами входных данных. Затем сценарий повторяется на всех интересующих исследователя компьютерах, в процессе чего происходит накопление экспериментов в общей базе данных.

\subsubsection{Логарифмический перебор}

В данном сценарии размер входных данных задаётся числом $N$ из заданного диапазона, причём количество строк и столбцов обрабатываемой матрицы принимаются одинаковыми. Таким образом, этот сценарий рассматривает только обработку квадратных матриц. Число $N$ меняется по закону $N = 2 ^ i$, где $i$ --- номер эксперимента.

\subsubsection{Случайный однообразный перебор}
В данном сценарии размер входных данных задаётся числом $N$ из заданного диапазона, означающим количество строк и столбцов в обрабатываемой матрице. Число генерируется случайным образом и распределено по диапазону равномерно. Этот сценарий рассматривает только обработку квадратных матриц. 

\subsubsection{Случайный перебор}
\label{random-exploration}

В данном сценарии размер входных данных задаётся парой чисел $N$ и $M$ из заданного диапазона, означающих количество строк и столбцов в обрабатываемой матрице, соответственно. Оба числа генерируются случайным образом и распределены по диапазону равномерно.

\subsection{Процесс работы системы в фоновом режиме}

В случае запуска системы на сервере возможна её работа в фоновом режиме. В таком случае алгоритм работы выглядит следующим образом.

\begin{enumerate}
\item Начало итерации. Пока не накоплено достаточное количество экспериментов в БД, ничего не делать. На этом шаге система находится в пассивном состоянии, ожидая ввода со стороны клиентов. Клиенты производят эксперименты в обычном режиме, при этом сохраняя их в удалённую БД, за которой следит сервер. Как только накоплено достаточное количество экспериментов, система переходит к следующему шагу.
\item Произвести обучение модели производительности программ на имеющихся экспериментах.
\item Произвести перекрёстную проверку модели на имеющихся данных. Если точность модели по результатам перекрёстной проверки ниже определённой константы, перейти к следующему шагу. Иначе, перейти к шагу 1 и закончить итерацию.
\item Произвести ранжирование признаков модели и отобрать на 1 признак больше, чем было включено в модель на предыдущей итерации. Перейти на шаг 2.
\end{enumerate}

Данный алгоритм обеспечивает периодическую подстройку модели под новые данные. Он является адаптивным и со временем начинает использовать всё больше и больше признаков экспериментов для увеличения точности модели. Его слабой стороной является требование наличия большого числа релевантных признаков, из которых можно выбирать на шаге 4. В данном случае система опирается на работу исследователей по добавлению модулей сбора всё новых и новых признаков.

При добавлении нового признака в модель изменение базы данных уже имеющихся экспериментов не требуется. Поскольку CouchDB является документо-ориентированной БД, она не требует наличия жёсткой схемы и новые поля могут добавляться в документы динамически. При этой наобходимо добавление кода в модуль связи с БД. Этот код должен отслеживать обращение к документам и в случае отсутствия необходимых полей (например, если получен документ старой версии, в котором ещё не было нового признака) производить действия по дополнению документов недостающими данными с помощью эмпирических правил.

\subsection{Предварительная обработка данных для анализа}

Система включает модуль предварительной обработки данных. Этот модуль связывается с БД, запрашивает документы, отвечающие определённому критерию, производит преобразование полученных документов в формат Comma Separated Value (CSV) и сохраняет их в файл.

\subsubsection{Критерии получения экспериментов из БД}

Для ограничения области изучения возможно получение из БД только тех документов, которые представляют интерес для исследователя в данный момент времени. Фильтрация документов возможна по следующим свойствам:

\begin{itemize}
\item дата и время проиведения эксперимента;
\item имя исследуемой программы;
\item название серии экспериментов (задаётся исследователем перед проведением серии экспериментов и используется для группировки).
\end{itemize}

Эти варианты доступа к экспериментам реализованы в виде представлений (\textit{англ. view}) в приложении CouchApp \cite{couchapp}. Представления программируются на языке JavaScript и представляют собой правила отбора документов по нужным критериям. Ключевой особенностью представлений CouchApp является то, что, будучи однажды построенным, представление производит очень быстрый отбор данных. В данном смысле оно подобно индексу в традиционно реляционной базе данных.

% Эксперименты
\section{Экспериментальная часть}

% Дальнейшая работа
\section{Заключение}
В рамках работы были решены задачи, перечисленные ниже.
\begin{itemize}
	\item Построена первоначальная образцовая модель производительности.
	\item Выявлены недостатки первоначальной модели.
	\item Разработана программа анализа модели производительности с целью выявления недостаточности имеющихся в модели признаков.
	\item Произведено добавление в модель признаков, позволяющих более точно описать производительность программы на данной платформе, в ручном и автоматическом режиме.
	\item Произведена оценка разработанной программы на наборе моделей, полученных с помощью анализа программ, входящих в состав набора Polybench.
\end{enumerate}
% Название библиографии
\renewcommand\bibname{Список использованных источников}
% Список использованных источников
\bibliography{biblio/course}
\end{document}
